\section{Recap Quiz}\label{sec:recap-quiz}

\begin{exercise}
  Let's say you have a game $\Game~G$ that you want to prove is hard, and you have the option to upper bound $\Pr[\game{G}{\adv} = 1]$ by $\advantage{AOMDL}{\bdv, \grgen}$, $\advantage{OMDL}{\cdv, \grgen}$, $\advantage{OMDL}{\ddv, \grgen}$ in the ROM\index{random oracle model rom@Random oracle model (ROM)}, or $\advantage{DL}{\edv, \grgen}$.
  What is your order of preference?
  \begin{enumerate}
    \item $\advantage{DL}{\edv, \grgen}$, $\advantage{AOMDL}{\bdv, \grgen}$, $\advantage{OMDL}{\cdv, \grgen}$, $\advantage{OMDL}{\ddv, \grgen}$ in the ROM
    \item $\advantage{DL}{\edv, \grgen}$, $\advantage{OMDL}{\cdv, \grgen}$, $\advantage{AOMDL}{\bdv, \grgen}$, $\advantage{OMDL}{\ddv, \grgen}$ in the ROM
    \item $\advantage{OMDL}{\ddv, \grgen}$ in the ROM, $\advantage{DL}{\edv, \grgen}$, $\advantage{AOMDL}{\bdv, \grgen}$, $\advantage{OMDL}{\cdv, \grgen}$
    \item $\advantage{AOMDL}{\bdv, \grgen}$, $\advantage{DL}{\edv, \grgen}$, $\advantage{OMDL}{\cdv, \grgen}$, $\advantage{OMDL}{\ddv, \grgen}$ in the ROM
  \end{enumerate}
\end{exercise}

\ifsolutions
\begin{mysolution}
  The answer is (1).
\end{mysolution}
\fi

\begin{exercise}
  Consider the following game hop\index{game@Game!game hopping@game hopping} from $\game{\Game~\algo{Guess}}{}$ to $\game{\Game~\algo{Guess}_2}{}$:
  
  \begin{figure}[h!]
    \begin{center}
      \begin{tcolorbox}[width=3cm]
        \begin{pchstack}[center]
            \procedure[headlinesep=1pt]{$\game{\Game~\algo{Guess}_2}{}$}{%
              n \defeq 2^\lambda\\
              x \sample \ZZ_n \\
              \gamechange{$\pcassert x \neq 0$} \\
              x' \gets \adv() \\
              \pcreturn x = x'
            }
        \end{pchstack}
      \end{tcolorbox}
    \end{center}
    \caption{Without the \gamechange{highlighted} change, this is the guessing game from \autoref{fig:guessing-game}. With the change, this is $\game{\Game~\algo{Guess}_2}{}$.}
  \end{figure}
  
  Which of the following is correct?
  \begin{enumerate}
    \item $\left|\advantage{Guess}{\adv} - \Pr[\game{\algo{Guess}_2}{} = 1]\right| \leq \Pr[x = 0]$
    \item $\left|\advantage{Guess}{\adv} - \Pr[\game{\algo{Guess}_2}{} = 1]\right| = \Pr[x = 0]$
    \item $\left|\advantage{Guess}{\adv} - \Pr[\game{\algo{Guess}_2}{} = 1]\right| \geq \Pr[x = 0]$
  \end{enumerate}
\end{exercise}

\ifsolutions
\begin{mysolution}
  The answer is (1).
  By the Difference Lemma\index{difference lemma@Difference lemma} (Lemma~\ref{lem:difference}), when the games are identical except when event $E$ occurs (here $E$ is ``$x = 0$''), the difference in success probabilities is bounded by $\Pr[E]$.
\end{mysolution}
\fi

\begin{exercise}
  Using the result from the previous exercise, what can we conclude about $\Pr[\game{\algo{Guess}_2}{} = 1]$?
  \begin{enumerate}
    \item $\Pr[\game{\algo{Guess}_2}{} = 1] \geq \advantage{Guess}{\adv} - 2^{-\lambda}$
    \item $\Pr[\game{\algo{Guess}_2}{} = 1] \geq \advantage{Guess}{\adv} + 2^{-\lambda}$
    \item $\Pr[\game{\algo{Guess}_2}{} = 1] < \advantage{Guess}{\adv} - 2^{-\lambda}$
  \end{enumerate}
\end{exercise}

\ifsolutions
\begin{mysolution}
  The answer is (1).
  
  From the previous exercise, we have $\left|\advantage{Guess}{\adv} - \Pr[\game{\algo{Guess}_2}{} = 1]\right| \leq 2^{-\lambda}$.
  
  This means:
  $-2^{-\lambda} \leq \advantage{Guess}{\adv} - \Pr[\game{\algo{Guess}_2}{} = 1] \leq 2^{-\lambda}$
  
  From the right inequality:
  $\advantage{Guess}{\adv} - \Pr[\game{\algo{Guess}_2}{} = 1] \leq 2^{-\lambda}$
  
  Rearranging:
  $\Pr[\game{\algo{Guess}_2}{} = 1] \geq \advantage{Guess}{\adv} - 2^{-\lambda}$
\end{mysolution}
\fi

\begin{exercise}
  Let $\adv$ be an algorithm with access to a random oracle\index{random oracle@Random oracle} $\pcoracle{H}: \{0,1\}^* \to \{0,1\}^\secpar$.
  The algorithm queries the random oracle exactly once with some $x \in \{0,1\}^\secpar$.
  What is $\Pr[\pcoracle{H}(x) = x]$?
  \begin{enumerate}
    \item $0$
    \item $2^{-\secpar}$
    \item $1$
    \item $2^{-2\secpar}$
  \end{enumerate}
\end{exercise}

\ifsolutions
\begin{mysolution}
  The answer is (2).
  The random oracle\index{random oracle@Random oracle} returns a uniformly random value from $\{0,1\}^\secpar$ on the first query to any input.
  Since there are $2^\secpar$ possible outputs and exactly one of them equals $x$, we have $\Pr[\pcoracle{H}(x) = x] = 2^{-\secpar}$.
\end{mysolution}
\fi

\begin{exercise}
  Let $p$ be a prime and let $\adv$ be an algorithm with access to a random oracle\index{random oracle@Random oracle} $\pcoracle{H}: \{0,1\}^* \to \ZZ_p$.
  The adversary queries the random oracle exactly twice with distinct $x_1, x_2 \in \ZZ_p$ (in that order) and receives $y_1 = \pcoracle{H}(x_1)$, $y_2 = \pcoracle{H}(x_2)$.
  What is $\Pr[y_2 = x_1 y_1 x_2^{-1}]$?
  \begin{enumerate}
    \item $0$
    \item $1/p$
    \item $1/(p-1)$
    \item $1$
  \end{enumerate}
\end{exercise}

\ifsolutions
\begin{mysolution}
  The answer is (2).
  
  Since $x_1 \neq x_2$, the queries are to different inputs, so $y_1$ and $y_2$ are independent uniform random values from $\ZZ_p$. 
  Given any fixed $x_1, x_2, y_1$, the value $x_1 y_1 x_2^{-1}$ is a fixed element of $\ZZ_p$.
  Since $y_2$ is uniformly random and independent of $y_1$, we have $\Pr[y_2 = x_1 y_1 x_2^{-1}] = 1/p$.
\end{mysolution}
\fi

\begin{exercise}
  Let $\gparam$ be a group description.
  Let $\adv$ be an algorithm with access to a random oracle\index{random oracle@Random oracle} $\pcoracle{H}: \{0,1\}^* \to \ZZ_p$ where $p$ is prime.
  The adversary queries the random oracle exactly twice with $(R_0, m_0)$ and $(R_1, m_1)$ where $R_i \in \GG$ and $m_i \in \{0,1\}^*$, receiving $y_0 = \pcoracle{H}(R_0, m_0)$ and $y_1 = \pcoracle{H}(R_1, m_1)$.
  What is $\Pr[R_0 \cdot g^{y_0} \cdot R_1^{-1} = g^{y_1}]$?
  \begin{enumerate}
    \item $0$
    \item $1/p$
    \item $1/(p-1)$
    \item $1$
  \end{enumerate}
\end{exercise}

\ifsolutions
\begin{mysolution}
  The answer is (2).
  
  Since the queries $(R_0, m_0)$ and $(R_1, m_1)$ are distinct, $y_0$ and $y_1$ are independent uniform random values from $\ZZ_p$.
  
  For any fixed $R_0, R_1 \in \GG$ and $y_0 \in \ZZ_p$, the value $R_0 \cdot g^{y_0} \cdot R_1^{-1}$ is a fixed element of $\GG$.
  Since $y_1$ is uniformly random in $\ZZ_p$, $g^{y_1}$ is uniformly random in $\GG$.
  Therefore, $\Pr[R_0 \cdot g^{y_0} \cdot R_1^{-1} = g^{y_1}] = 1/p$.
  
  Note: If this event occurred, we would have $R_0 \cdot g^{\pcoracle{H}(R_0, m_0)} = R_1 \cdot g^{\pcoracle{H}(R_1, m_1)}$, meaning
  \[
    \algo{TapCom.Commit}(R_0, m_0) = \algo{TapCom.Commit}(R_1, m_1)
  \] and the adversary would have won $\game{\Game~\algo{ComBind}}{\algo{TapCom}}$ in the random oracle\index{random oracle@Random oracle} model\index{random oracle model rom@Random oracle model (ROM)}.
  In \autoref{lem:taproot-binding}, we generalize this to $q$ random oracle queries.
\end{mysolution}
\fi

\begin{exercise}
  The following demonstrates how games in the random oracle\index{random oracle@Random oracle} model\index{random oracle model rom@Random oracle model (ROM)} can examine the table $T$ of previous queries.
  Let $\algo{H}$ be a random oracle mapping to $\{0,1\}^\secpar$ and let $\adv$ be an algorithm making at most $q_h$ queries to $\algo{H}$.
  Let $\game{\Game~G}{}$ be the game defined in Figure~\ref{fig:quiz-ro-game}.
  What is $\Pr[\game{G}{} = 1]$?
  \begin{enumerate}
    \item $0$
    \item $1$
    \item $\leq \frac{q_h}{2^\secpar}$
    \item $\leq \frac{(q_h + 1) q_h}{2^\secpar}$
  \end{enumerate}
\end{exercise}

\begin{figure}[tbhp]
  \begin{center}
    \begin{tcolorbox}[width=0.6\textwidth]
      \begin{pchstack}[center]
        \procedure[headlinesep=1pt]{$\game{\Game~G}{}$}{%
          T \defeq \emptyset \\
          (x, h) \gets \adv^{\pcoracle{H}}() \\
          \pcif \pcoracle{H}(x) \neq h \pcthen \\
          \t \pcreturn \pcfalse \\
          \pcif \exists x' \neq x : T[x'] = h \pcthen \\
          \t \pcreturn \pctrue \\
          \pcreturn \pcfalse
        }
        \procedure[headlinesep=1pt]{Oracle $\pcoracle{H}(x)$}{%
          \pcif T[x] = \bot \pcthen \\
          \t T[x] \sample \{0,1\}^\secpar \\
          \pcreturn T[x]
        }
      \end{pchstack}
    \end{tcolorbox}
  \end{center}
  \caption{Game $\game{\Game~G}{}$ for random oracle quiz question}
  \label{fig:quiz-ro-game}
\end{figure}

\ifsolutions
\begin{mysolution}
  The answer is (3).
  
  For the game to return $\pctrue$, two conditions must be satisfied:
  \begin{enumerate}
    \item $\pcoracle{H}(x) = h$ (the adversary's predicted hash value must be correct)
    \item $\exists x' \neq x : T[x'] = h$ (there must be a collision: a different input $x'$ that was previously queried with the same hash $h$)
  \end{enumerate}
  
  Since both conditions must hold:
  \[\Pr[\game{G}{} = 1] \leq \Pr[\exists x' \neq x : \pcoracle{H}(x') = \pcoracle{H}(x)]\]
  
  The adversary makes at most $q_h$ queries during its execution, and the game makes one additional query $\pcoracle{H}(x)$.
  So there are at most $q_h + 1$ total queries.
  For each query, the probability of colliding with any of the at most $q_h$ entries already in the table is at most $\frac{q_h}{2^\secpar}$.
  
  By the union bound\index{union bound@Union bound} over all $q_h + 1$ queries:
  \[
    \Pr[\exists x' \neq x : \pcoracle{H}(x') = \pcoracle{H}(x)] \leq (q_h + 1) \cdot \frac{q_h}{2^\secpar}
  \]
  
  Therefore, $\Pr[\game{G}{} = 1] \leq \frac{(q_h + 1)q_h}{2^\secpar}$.
\end{mysolution}
\fi

\begin{exercise}
  Which of the following represents a gap between security proofs and real-world security?
  \begin{enumerate}
    \item In practice, $\lambda$ is a fixed value (e.g., 128 or 256)
    \item Concrete hash functions\index{hash function@Hash function} like SHA256 are not hash functions
    \item Unfalsifiable\index{falsifiable assumption@Falsifiable assumption} assumptions cannot be validated experimentally
    \item Security models may not capture all attack vectors
    \item Random oracles\index{random oracle@Random oracle} do not exist
    \item Security proofs don't model users writing keys on post-it notes
    \item All of the above
  \end{enumerate}
\end{exercise}

\ifsolutions
\begin{mysolution}
  The answer is (7).
\end{mysolution}
\fi
